\documentclass[11pt]{article}
\usepackage[utf8]{inputenc}
\usepackage{geometry}
\geometry{a4paper, margin=1in}
\usepackage{hyperref}
\usepackage{booktabs}
\usepackage{enumitem}
\usepackage{float}
\usepackage{xcolor}
\usepackage{fancyvrb}

\title{Landslide Identification Using Machine Learning \\ \large Run Guide --- Clone to Output}
\author{}
\date{}

\begin{document}

\maketitle
\tableofcontents
\newpage

% ============================================================================
\section{All Commands --- Clone to Output (Copy-Paste)}
% ============================================================================

Run these commands one by one in your terminal. This is the complete sequence from cloning the repository to seeing the prediction output.

\begin{verbatim}
# 1. Clone the repository
git clone https://github.com/<your-username>/Landslide_Identification_Using_ML.git
cd Landslide_Identification_Using_ML

# 2. Create conda environment
conda create -n landslide-ml python=3.10 -y
conda activate landslide-ml

# 3. Install PyTorch (pick ONE based on your GPU)
#    Option A: RTX 20/30/40 series
pip install torch torchvision --index-url https://download.pytorch.org/whl/cu124
#    Option B: RTX 50-series (Blackwell)
pip install --pre torch torchvision --index-url https://download.pytorch.org/whl/nightly/cu128
#    Option C: No GPU (CPU only)
pip install torch torchvision

# 4. Install all other dependencies
cd project
pip install -r requirements.txt

# 5. Verify everything is installed
python -c "import torch; print('PyTorch:', torch.__version__); print('CUDA:', torch.cuda.is_available())"
python -c "import hmmlearn; print('hmmlearn: OK')"

# 6. Train both models (skip if checkpoints are already in checkpoints/)
python pipeline/run_pipeline.py --mode train

# 7. Predict on a landslide image
python pipeline/run_pipeline.py --mode predict \
    --image data/processed/test/debris_flow/ls_test_0011.jpg \
    --country India --forecast 3

# 8. Predict on a non-landslide image
python pipeline/run_pipeline.py --mode predict \
    --image data/processed/test/non_landslide/nls_test_0001.jpg \
    --country India

# 9. Evaluate model performance (prints accuracy, F1, AUC + saves plots)
python phase1_alexnet/evaluate.py

# 10. Run HMM predictions for multiple countries
python phase2_hmm/hmm_predict.py

# 11. Run full pipeline evaluation
python pipeline/run_pipeline.py --mode evaluate
\end{verbatim}

\noindent \textbf{On Windows}, replace \texttt{\textbackslash} with a single line:
\begin{verbatim}
python pipeline/run_pipeline.py --mode predict --image data\processed\test\debris_flow\ls_test_0011.jpg --country India --forecast 3
\end{verbatim}

\noindent After running these commands:
\begin{itemize}[noitemsep]
    \item Commands 7--8 print the landslide prediction with confidence, type, and forecast
    \item Command 9 prints accuracy, precision, recall, F1, ROC-AUC and saves \texttt{plots/confusion\_matrix.png} and \texttt{plots/roc\_curve.png}
    \item Command 10 prints country-wise risk profiles with peak risk months
    \item Command 11 prints overall pipeline evaluation metrics
\end{itemize}

\newpage

% ============================================================================
\section{What This Project Does}
% ============================================================================

This is a two-phase ML pipeline that takes a satellite image as input and produces a complete landslide risk report as output.

\begin{itemize}
    \item \textbf{Phase 1 (Image Classification):} An ensemble of EfficientNet-B3 + ViT-B/16 classifies the image as landslide or non-landslide with a confidence score.
    \item \textbf{Phase 2 (Risk Forecasting):} A Hidden Markov Model trained on 9,471 real NASA landslide events predicts the landslide type, occurrence probability, peak risk month, and a multi-step future forecast.
\end{itemize}

% ============================================================================
\section{Step 1 --- Clone the Repository}
% ============================================================================

\begin{verbatim}
git clone https://github.com/<your-username>/Landslide_Identification_Using_ML.git
cd Landslide_Identification_Using_ML
\end{verbatim}

\noindent After cloning, your directory structure should look like:

\begin{verbatim}
Landslide_Identification_Using_ML/
  project/
    checkpoints/          <-- pre-trained model weights
    data/
      processed/          <-- satellite images (train/val/test)
      excel/              <-- NASA Global Landslide Catalog
    phase1_alexnet/       <-- CNN training, evaluation, prediction
    phase2_hmm/           <-- HMM training and prediction
    pipeline/             <-- end-to-end pipeline script
    notebooks/            <-- Jupyter notebooks (step-by-step)
    plots/                <-- generated plots (confusion matrix, ROC, Grad-CAM)
    utils/                <-- helpers (Grad-CAM, plotting, focal loss)
    config.py             <-- all hyperparameters and paths
    requirements.txt      <-- Python dependencies
\end{verbatim}

% ============================================================================
\section{Step 2 --- Set Up the Environment}
% ============================================================================

\subsection{2.1 --- Create a conda environment}

\begin{verbatim}
conda create -n landslide-ml python=3.10 -y
conda activate landslide-ml
\end{verbatim}

\subsection{2.2 --- Install PyTorch (GPU)}

\textbf{For most NVIDIA GPUs (RTX 20/30/40 series):}
\begin{verbatim}
pip install torch torchvision --index-url https://download.pytorch.org/whl/cu124
\end{verbatim}

\textbf{For RTX 50-series GPUs (Blackwell / sm\_120):}
\begin{verbatim}
pip install --pre torch torchvision \
    --index-url https://download.pytorch.org/whl/nightly/cu128
\end{verbatim}

\textbf{CPU only (no GPU):}
\begin{verbatim}
pip install torch torchvision
\end{verbatim}

\subsection{2.3 --- Install remaining dependencies}

\begin{verbatim}
cd project
pip install -r requirements.txt
\end{verbatim}

\subsection{2.4 --- Verify installation}

\begin{verbatim}
python -c "import torch; print('PyTorch:', torch.__version__)"
python -c "import torch; print('CUDA available:', torch.cuda.is_available())"
python -c "import hmmlearn; print('hmmlearn OK')"
\end{verbatim}

All three commands should run without errors. CUDA should show \texttt{True} if you have a GPU.

% ============================================================================
\section{Step 3 --- Navigate to the Project Directory}
% ============================================================================

\textbf{All commands from this point must be run from inside the \texttt{project/} directory.}

\begin{verbatim}
cd project
\end{verbatim}

Or the full path (adjust to your system):
\begin{verbatim}
cd Landslide_Identification_Using_ML/project
\end{verbatim}

% ============================================================================
\section{Step 4 --- Quick Demo (Fastest Way to See Output)}
% ============================================================================

If the pre-trained checkpoints are already included in the repo, you can skip training and go straight to prediction.

\subsection{4.1 --- Predict on a single image}

\textbf{Linux / macOS:}
\begin{verbatim}
python pipeline/run_pipeline.py \
    --mode predict \
    --image data/processed/test/debris_flow/ls_test_0011.jpg \
    --country India \
    --forecast 3
\end{verbatim}

\textbf{Windows (single line):}
\begin{verbatim}
python pipeline/run_pipeline.py --mode predict --image
    data\processed\test\debris_flow\ls_test_0011.jpg --country India
    --forecast 3
\end{verbatim}

\subsection{4.2 --- Expected output}

If a landslide is detected:
\begin{verbatim}
Phase 1 - LANDSLIDE (confidence: 95.9%)
Phase 2 - Type: Landslide
          Occurrence probability: 30.0%
          Peak risk month: July
          Future forecast:
            Step 1: Landslide (prob=26.4%)
            Step 2: Landslide (prob=21.4%)
            Step 3: Landslide (prob=18.1%)

LANDSLIDE DETECTED - Type: Landslide (occurrence probability: 30%)
\end{verbatim}

If no landslide is detected:
\begin{verbatim}
NO LANDSLIDE DETECTED (confidence: 78%)
\end{verbatim}

\subsection{4.3 --- Try different images}

Test with a landslide image:
\begin{verbatim}
python pipeline/run_pipeline.py --mode predict \
    --image data/processed/test/debris_flow/ls_test_0015.jpg \
    --country Nepal
\end{verbatim}

Test with a non-landslide image:
\begin{verbatim}
python pipeline/run_pipeline.py --mode predict \
    --image data/processed/test/non_landslide/nls_test_0001.jpg \
    --country India
\end{verbatim}

\subsection{4.4 --- Run HMM predictions for multiple countries}

\begin{verbatim}
python phase2_hmm/hmm_predict.py
\end{verbatim}

This shows type, occurrence probability, peak risk month, historical statistics, and 3-step forecasts for India, Nepal, Philippines, and United States.

% ============================================================================
\section{Step 5 --- Training from Scratch (Optional)}
% ============================================================================

If checkpoints are not included or you want to retrain:

\subsection{5.1 --- Train everything at once}

\begin{verbatim}
python pipeline/run_pipeline.py --mode train
\end{verbatim}

This trains Phase 1 (CNN, 30 epochs) and Phase 2 (HMM, 100 iterations). On a GPU this takes approximately 20--30 minutes. On CPU it takes 4--6 hours.

\subsection{5.2 --- Train models separately}

\textbf{Phase 1 --- CNN (EfficientNet-B3):}
\begin{verbatim}
python phase1_alexnet/train.py
\end{verbatim}

\textbf{Phase 2 --- HMM:}
\begin{verbatim}
python phase2_hmm/hmm_model.py
\end{verbatim}

Trained models are saved to \texttt{checkpoints/} automatically.

% ============================================================================
\section{Step 6 --- Evaluate and Generate Plots}
% ============================================================================

\subsection{6.1 --- Evaluate Phase 1 (CNN metrics)}

\begin{verbatim}
python phase1_alexnet/evaluate.py
\end{verbatim}

\textbf{What it prints:} Accuracy, Precision, Recall, F1-Score, ROC-AUC on 699 test images.

\textbf{What it saves:}
\begin{itemize}[noitemsep]
    \item \texttt{plots/confusion\_matrix.png} --- confusion matrix heatmap
    \item \texttt{plots/roc\_curve.png} --- ROC curve with AUC value
\end{itemize}

\subsection{6.2 --- Evaluate full pipeline}

\begin{verbatim}
python pipeline/run_pipeline.py --mode evaluate
\end{verbatim}

Runs Phase 1 evaluation on the full test set and reports HMM log-likelihood.

\subsection{6.3 --- Generated plots for presentation}

After evaluation, the \texttt{plots/} folder contains:

\begin{table}[H]
    \centering
    \renewcommand{\arraystretch}{1.2}
    \begin{tabular}{lp{9cm}}
        \toprule
        \textbf{File} & \textbf{Description} \\
        \midrule
        \texttt{confusion\_matrix.png} & Shows true vs predicted labels for all test images \\
        \texttt{roc\_curve.png} & ROC curve with AUC score \\
        \texttt{training\_history.png} & Loss and accuracy over training epochs \\
        \texttt{gradcam\_ls\_test\_0011.png} & Grad-CAM heatmap showing which image regions the model focused on \\
        \bottomrule
    \end{tabular}
\end{table}

% ============================================================================
\section{Step 7 --- Running Jupyter Notebooks (Alternative)}
% ============================================================================

For a step-by-step walkthrough with inline outputs:

\begin{verbatim}
pip install jupyter
jupyter notebook notebooks/
\end{verbatim}

Run the notebooks in order:

\begin{table}[H]
    \centering
    \renewcommand{\arraystretch}{1.3}
    \begin{tabular}{clp{8cm}}
        \toprule
        \textbf{\#} & \textbf{Notebook} & \textbf{What It Does} \\
        \midrule
        0 & 00\_setup\_and\_data\_download & Verifies dataset, displays sample images \\
        1 & 01\_phase1\_alexnet\_training & Trains CNN (30 epochs), shows loss curves \\
        2 & 02\_phase1\_evaluation & Evaluates test set, prints metrics, saves plots \\
        3 & 03\_phase2\_hmm\_training & Trains HMM on NASA catalog \\
        4 & 04\_phase2\_hmm\_evaluation & Shows transition matrix, country predictions \\
        5 & 05\_full\_pipeline\_demo & End-to-end: image $\to$ Phase 1 $\to$ Phase 2 $\to$ risk report \\
        \bottomrule
    \end{tabular}
\end{table}

% ============================================================================
\section{Project File Reference}
% ============================================================================

\subsection{Key configuration (\texttt{config.py})}

\begin{table}[H]
    \centering
    \renewcommand{\arraystretch}{1.2}
    \begin{tabular}{lll}
        \toprule
        \textbf{Parameter} & \textbf{Value} & \textbf{Meaning} \\
        \midrule
        NUM\_EPOCHS & 30 & Training epochs for CNN \\
        BATCH\_SIZE & 32 & Images per training batch \\
        NUM\_CLASSES & 6 & Classification categories \\
        PHASE1\_THRESHOLD & 0.467 & Optimal F1 threshold for landslide detection \\
        HMM\_N\_COMPONENTS & 8 & Hidden states in HMM \\
        ENSEMBLE\_WEIGHT\_EFFNET & 0.6 & EfficientNet-B3 weight in ensemble (60\%) \\
        ENSEMBLE\_WEIGHT\_ALEXNET & 0.4 & ViT-B/16 weight in ensemble (40\%) \\
        \bottomrule
    \end{tabular}
\end{table}

\subsection{Pre-trained model files (\texttt{checkpoints/})}

\begin{table}[H]
    \centering
    \renewcommand{\arraystretch}{1.2}
    \begin{tabular}{lp{9cm}}
        \toprule
        \textbf{File} & \textbf{Description} \\
        \midrule
        \texttt{efficientnet\_b3\_best.pth} & EfficientNet-B3 CNN weights (Phase 1) \\
        \texttt{vit\_b\_16\_best.pth} & Vision Transformer weights (Phase 1) \\
        \texttt{hmm\_model.pkl} & Trained HMM with 8 hidden states, 42 symbols (Phase 2) \\
        \texttt{hmm\_encoder.pkl} & LabelEncoder for landslide types \\
        \texttt{hmm\_trigger\_encoder.pkl} & LabelEncoder for trigger types \\
        \bottomrule
    \end{tabular}
\end{table}

These are loaded automatically during prediction. No retraining needed unless desired.

\subsection{Dataset (\texttt{data/})}

\begin{table}[H]
    \centering
    \renewcommand{\arraystretch}{1.2}
    \begin{tabular}{llcc}
        \toprule
        \textbf{Split} & \textbf{Location} & \textbf{Landslide} & \textbf{Non-Landslide} \\
        \midrule
        Train & \texttt{data/processed/train/} & 616 & 1,574 \\
        Validation & \texttt{data/processed/val/} & 158 & 392 \\
        Test & \texttt{data/processed/test/} & 211 & 488 \\
        \midrule
        Phase 2 (HMM) & \texttt{data/excel/nasa\_glc.csv} & \multicolumn{2}{c}{9,471 events, 141 countries} \\
        \bottomrule
    \end{tabular}
\end{table}

% ============================================================================
\section{Troubleshooting}
% ============================================================================

\begin{table}[H]
    \centering
    \renewcommand{\arraystretch}{1.5}
    \begin{tabular}{p{5.5cm}p{8.5cm}}
        \toprule
        \textbf{Problem} & \textbf{Solution} \\
        \midrule
        \texttt{ModuleNotFoundError} & Make sure you activated the environment: \texttt{conda activate landslide-ml} \\
        \texttt{FileNotFoundError} for image & Use the correct path relative to \texttt{project/} directory \\
        \texttt{conda} command not found & \textbf{Linux/macOS:} \texttt{export PATH="\$HOME/miniconda3/bin:\$PATH"} \newline \textbf{Windows PowerShell:} Run \texttt{conda init powershell} then restart terminal \\
        GPU not detected & Check: \texttt{python -c "import torch; print(torch.cuda.is\_available())"} \newline If False, reinstall PyTorch with CUDA (see Step 2.2) \\
        RTX 50-series GPU error \newline (\texttt{sm\_120 not compatible}) & Install PyTorch nightly with cu128 (see Step 2.2) \\
        Checkpoint not found & Train first: \texttt{python pipeline/run\_pipeline.py --mode train} \\
        HMM shows 0 records & Verify \texttt{data/excel/} contains the NASA GLC CSV file \\
        \texttt{CUDA out of memory} & Reduce \texttt{BATCH\_SIZE} in \texttt{config.py} (try 16 or 8) \\
        \bottomrule
    \end{tabular}
\end{table}

% ============================================================================
\section{Quick Reference --- Commands Summary}
% ============================================================================

\begin{table}[H]
    \centering
    \renewcommand{\arraystretch}{1.4}
    \begin{tabular}{p{4.5cm}p{9.5cm}}
        \toprule
        \textbf{Task} & \textbf{Command} \\
        \midrule
        Setup environment & \texttt{conda create -n landslide-ml python=3.10 -y} \\
        Install dependencies & \texttt{pip install -r requirements.txt} \\
        Train everything & \texttt{python pipeline/run\_pipeline.py --mode train} \\
        Predict single image & \texttt{python pipeline/run\_pipeline.py --mode predict --image <path> --country India} \\
        Evaluate CNN & \texttt{python phase1\_alexnet/evaluate.py} \\
        Evaluate full pipeline & \texttt{python pipeline/run\_pipeline.py --mode evaluate} \\
        HMM country predictions & \texttt{python phase2\_hmm/hmm\_predict.py} \\
        Launch notebooks & \texttt{jupyter notebook notebooks/} \\
        \bottomrule
    \end{tabular}
\end{table}

\end{document}
